% ============================================================
% 第5章 习题提示
% ============================================================
\section{第5章习题提示}

需要提示时再查阅。先尝试独立完成，卡住后逐级查看。

\subsection*{练习5.14（难度3·普通）}

\textbf{L2 提醒：} 已知 \(\lg 2\) 和 \(\lg 3\)，把 6、12、1.5 分解成 2 和 3 的乘积或商。
\(6 = 2 \times 3\)，\(12 = 2^2 \times 3\)，\(1.5 = 3/2\)。然后用对数运算法则展开。

\subsection*{练习5.15（难度3·普通）}

\textbf{L2 提醒：}
用对数运算法则合并左边两项：\(\lg(x+2) + \lg(x-2) = \lg[(x+2)(x-2)]\)。
然后两边同时去掉 \(\lg\)，解普通方程。注意验根！

\subsection*{练习5.16（难度3·普通）}

\textbf{L2 提醒：}
(1) 把 \(M=5\) 代入 \(\lg E = 4.8 + 1.5M\) 求 \(E\)。
(2) 分别算震级 5 和 6 的能量，后者除以前者。
注意：\(10^{\lg E} = E\)（对数的定义）。

\subsection*{练习5.17（难度3·普通）}

\textbf{L2 提醒：}
用换底公式写成：
\[
\frac{\lg 3}{\lg 2} \times \frac{\lg 4}{\lg 3} \times \frac{\lg 5}{\lg 4}
\]
分子分母逐项约分，最后剩下什么？

\subsection*{练习5.18（难度3·普通）}

\textbf{L2 提醒：}
用换底公式 \(\log_a b = \dfrac{\lg b}{\lg a}\)，那么 \(\log_b a = \dfrac{\lg a}{\lg b}\)。
两者相乘看看结果。

\subsection*{练习5.19（难度4·中等）}

\textbf{L3 第一步：}
(1) 用对数运算法则合并左边：
\(\log_2(x+1) + \log_2(x-1) = \log_2[(x+1)(x-1)] = \log_2(x^2 - 1) = 3\)。
(2) 令 \(t = \lg x\)，原方程化为 \(t^2 - t - 2 = 0\)。

\textbf{L4 完整思路：}
(1) \(\log_2(x^2 - 1) = 3 \Rightarrow x^2 - 1 = 2^3 = 8 \Rightarrow x^2 = 9 \Rightarrow x = \pm 3\)。
验根：\(x=3\) 时 \(x+1>0, x-1>0\) [OK]；\(x=-3\) 时 \(x+1<0\) [X]。所以 \(x=3\)。
(2) \(t^2 - t - 2 = 0 \Rightarrow (t-2)(t+1)=0 \Rightarrow t=2\) 或 \(t=-1\)。
\(\lg x = 2 \Rightarrow x = 100\)；\(\lg x = -1 \Rightarrow x = 0.1\)。验根均通过。

\subsection*{练习5.20（难度4·中等）}

\textbf{L3 第一步：}
(1) \(f(0) = \ln(0 + \sqrt{1+0}) = \ln 1 = 0\)。
(2) 奇偶性判断：计算 \(f(-x)\)。

\textbf{L4 完整思路：}
(1) \(f(0)=0\)，\(f(1)=\ln(1+\sqrt{2})\)，\(f(-1)=\ln(-1+\sqrt{2}) = \ln(\sqrt{2}-1)\)。
(2) \(f(-x) = \ln(-x + \sqrt{1+x^2})\)。注意到 \(\sqrt{1+x^2} - x = \frac{1}{\sqrt{1+x^2}+x}\)（有理化）。
\(\therefore\) \(f(-x) = \ln\frac{1}{\sqrt{1+x^2}+x} = -\ln(x+\sqrt{1+x^2}) = -f(x)\)。奇函数。

\subsection*{练习5.21（难度4·中等）}

\textbf{L3 第一步：}
把已知的 \(a = \log_2 3\) 和 \(b = \log_3 5\) 用换底公式写成常用对数形式：
\(a = \frac{\lg 3}{\lg 2}\)，\(b = \frac{\lg 5}{\lg 3}\)。然后 \(\log_2 15 = \frac{\lg 15}{\lg 2}\)。

\textbf{L4 完整思路：}
\(\log_2 15 = \frac{\lg 15}{\lg 2} = \frac{\lg(3 \times 5)}{\lg 2} = \frac{\lg 3 + \lg 5}{\lg 2}\)。
\(a = \frac{\lg 3}{\lg 2} \Rightarrow \lg 3 = a \lg 2\)；\(b = \frac{\lg 5}{\lg 3} \Rightarrow \lg 5 = b \lg 3 = ab \lg 2\)。
代入得 \(\frac{a\lg 2 + ab\lg 2}{\lg 2} = a + ab = a(1+b)\)。

\subsection*{练习5.22（难度4·中等）}

\textbf{L3 第一步：}
平移变换：\(y = \log_2(x-1)\) 是 \(y = \log_2 x\) 向右平移 1 个单位；
再加 2 是向上平移 2 个单位。定义域由 \(x-1 > 0\) 确定。值域由平移量确定。

\textbf{L4 完整思路：}
\(y = \log_2 x\) → 右移 1 → \(\log_2(x-1)\) → 上移 2 → \(\log_2(x-1)+2\)。
定义域：\(x-1 > 0 \Rightarrow x > 1\)，即 \((1, +\infty)\)。
值域：\(\log_2(x-1)\) 值域 \(\mathbb{R}\)，加 2 后仍是 \(\mathbb{R}\)。

\subsection*{练习5.23（难度4·中等）}

\textbf{L3 第一步：}
(1) 底 2 > 1，对数函数递增，所以 \(\log_2 x > 3 = \log_2 2^3\)。
(2) 底 0.5 < 1，对数函数递减，所以 \(\log_{0.5} x \ge -2 = \log_{0.5} (0.5)^{-2}\)。

\textbf{L4 完整思路：}
(1) \(\log_2 x > \log_2 8 \Rightarrow x > 8\)。
(2) \(\log_{0.5} x \ge \log_{0.5} 4 \Rightarrow x \le 4\)（注意底<1 时不等号反向）。
结合定义域 \(x > 0\)，得 \(0 < x \le 4\)。

\subsection*{练习5.24（难度5·进阶）}

\textbf{L3 第一步：}
先把右边拆开：\(\log_2(2^{x+1} - 3)\) 不能直接化简。但注意 \(x + \log_2(2^{x+1} - 3) = \log_2(2^x) + \log_2(2^{x+1} - 3) = \log_2[2^x(2^{x+1} - 3)]\)。
左边不动：\(\log_2(4^x + 4)\)。去掉 \(\log_2\) 得到方程。

\textbf{L4 完整思路：}
\(4^x + 4 = 2^x(2^{x+1} - 3)\)。令 \(t = 2^x\)，则 \(4^x = t^2\)。
\(t^2 + 4 = t(2t - 3) \Rightarrow t^2 + 4 = 2t^2 - 3t \Rightarrow t^2 - 3t - 4 = 0\)。
\((t-4)(t+1)=0 \Rightarrow t=4\)（\(t=-1\) 舍）。\(2^x = 4 \Rightarrow x = 2\)。

\subsection*{练习5.25（难度5·进阶）}

\textbf{L3 第一步：}
(1) 奇偶性：计算 \(f(-x) = \ln\frac{1+x}{1-x}\)。
(2) 把 \(f(a) + f(b)\) 写成 \(\ln\) 的形式，利用对数运算法则合并。

\textbf{L4 完整思路：}
(1) \(f(-x) = \ln\frac{1+x}{1-x} = -\ln\frac{1-x}{1+x} = -f(x)\)，奇函数。
(2) \(f(a) + f(b) = \ln\frac{1-a}{1+a} + \ln\frac{1-b}{1+b} = \ln\left[\frac{(1-a)(1-b)}{(1+a)(1+b)}\right]\)。
化简后者：\(\frac{1 - \frac{a+b}{1+ab}}{1 + \frac{a+b}{1+ab}} = \frac{(1+ab)-(a+b)}{(1+ab)+(a+b)} = \frac{(1-a)(1-b)}{(1+a)(1+b)}\)。
所以 \(f(a) + f(b) = f\left(\frac{a+b}{1+ab}\right)\)。

\subsection*{练习5.26（难度5·进阶）}

\textbf{L3 第一步：}
(1) 代入 \(I = 10^{-6}\) 到公式 \(d = 10\lg\frac{I}{I_0}\)，其中 \(I_0 = 10^{-12}\)。
(2) 代入 \(I = 1\)。
(3) 比较 \(d+10\) 和 \(d\) 对应的强度。

\textbf{L4 完整思路：}
(1) \(d = 10\lg\frac{10^{-6}}{10^{-12}} = 10\lg 10^6 = 10 \times 6 = 60\) 分贝。
(2) \(d = 10\lg\frac{1}{10^{-12}} = 10\lg 10^{12} = 10 \times 12 = 120\) 分贝。
(3) 设原强度 \(I\) 对应分贝 \(d\)，新强度 \(I'\) 对应 \(d+10\)：
\(d+10 = 10\lg\frac{I'}{I_0} \Rightarrow 10\lg\frac{I'}{I_0} - 10\lg\frac{I}{I_0} = 10 \Rightarrow \lg\frac{I'}{I} = 1 \Rightarrow I' = 10I\)。
分贝每增加 10，强度扩大 10 倍。

\subsection*{练习5.27（难度5·进阶）}

\textbf{L3 第一步：}
分类讨论底数 \(a > 1\) 和 \(0 < a < 1\) 两种情况。
分别比较 \(a^2 + 1\) 和 \(2a\) 的大小，再结合对数函数的单调性。

\textbf{L4 完整思路：}
先比较 \(a^2 + 1\) 和 \(2a\) 的大小：\(a^2 + 1 - 2a = (a-1)^2 \ge 0\)，所以 \(a^2 + 1 \ge 2a\)。
当 \(a > 1\) 时，\(\log_a x\) 递增，\(\log_a(a^2+1) \ge \log_a(2a)\)。
当 \(0 < a < 1\) 时，\(\log_a x\) 递减，\(\log_a(a^2+1) \le \log_a(2a)\)。
等号在 \(a = 1\) 时取到，但 \(a \neq 1\)，所以严格不等。

\subsection*{练习5.28（难度6·拔高）}

\textbf{L2 知识点梳理：}
(1) 定义域为 \(\mathbb{R}\) 意味着 \(x^2 - 2ax + 3 > 0\) 对一切实数 \(x\) 成立 → 判别式 \(\Delta < 0\)。
(2) 值域为 \(\mathbb{R}\) 意味着 \(x^2 - 2ax + 3\) 必须能取到所有正数 → 最小值 \(\le 0\) → 判别式 \(\Delta \ge 0\)。

\textbf{L3 第一步：}
令 \(g(x) = x^2 - 2ax + 3\)，对称轴 \(x = a\)，最小值 \(g(a) = a^2 - 2a^2 + 3 = 3 - a^2\)。
(1) 定义域 \(\mathbb{R}\) ⟺ \(g(x) > 0\) 恒成立 ⟺ 最小值 \(> 0\)。
(2) 值域 \(\mathbb{R}\) ⟺ \(g(x)\) 能取遍所有正数 ⟺ 最小值 \(\le 0\)。

\textbf{L4 完整思路：}
(1) \(3 - a^2 > 0 \Rightarrow a^2 < 3 \Rightarrow -\sqrt{3} < a < \sqrt{3}\)。
(2) \(3 - a^2 \le 0 \Rightarrow a^2 \ge 3 \Rightarrow a \le -\sqrt{3}\) 或 \(a \ge \sqrt{3}\)。

\subsection*{练习5.29（难度6·拔高）}

\textbf{L2 知识点梳理：}
方程两边底数不同（2 和 4），需要用换底公式统一底数。
观察到 \(4 = 2^2\)，所以 \(\log_4(x+5) = \frac{\log_2(x+5)}{\log_2 4} = \frac{\log_2(x+5)}{2}\)。

\textbf{L3 第一步：}
\(\log_2(x+1) = \frac{\log_2(x+5)}{2} \Rightarrow 2\log_2(x+1) = \log_2(x+5)\)。
\(\log_2(x+1)^2 = \log_2(x+5)\)。

\textbf{L4 完整思路：}
\((x+1)^2 = x+5 \Rightarrow x^2 + 2x + 1 = x + 5 \Rightarrow x^2 + x - 4 = 0\)。
\(x = \frac{-1 \pm \sqrt{1+16}}{2} = \frac{-1 \pm \sqrt{17}}{2}\)。
验根：\(x+1 > 0\) 且 \(x+5 > 0\)。\(\frac{-1+\sqrt{17}}{2} > 0\) [OK]；\(\frac{-1-\sqrt{17}}{2} < 0\)，代入 \(x+1 < 0\) [X]。
所以 \(x = \frac{-1 + \sqrt{17}}{2}\)。

\subsection*{练习5.30（难度6·拔高）}

\textbf{L2 知识点梳理：}
已知 \(\log_{12} 27 = a\)，求 \(\log_6 16\)。思路：把两者都写成常用对数形式，找到 \(\lg 2\)、\(\lg 3\) 的关系。

\textbf{L3 第一步：}
\(\log_{12} 27 = \frac{\lg 27}{\lg 12} = \frac{\lg 3^3}{\lg(2^2 \times 3)} = \frac{3\lg 3}{2\lg 2 + \lg 3} = a\)。
由此解出 \(\frac{\lg 3}{\lg 2}\) 的关系。
\(\log_6 16 = \frac{\lg 16}{\lg 6} = \frac{\lg 2^4}{\lg(2 \times 3)} = \frac{4\lg 2}{\lg 2 + \lg 3}\)。

\textbf{L4 完整思路：}
由 \(\frac{3\lg 3}{2\lg 2 + \lg 3} = a \Rightarrow 3\lg 3 = a(2\lg 2 + \lg 3) \Rightarrow (3-a)\lg 3 = 2a\lg 2\)。
\(\frac{\lg 3}{\lg 2} = \frac{2a}{3-a}\)。
\(\log_6 16 = \frac{4\lg 2}{\lg 2 + \lg 3} = \frac{4}{\frac{\lg 2 + \lg 3}{\lg 2}} = \frac{4}{1 + \frac{\lg 3}{\lg 2}} = \frac{4}{1 + \frac{2a}{3-a}} = \frac{4(3-a)}{3-a+2a} = \frac{4(3-a)}{3+a}\)。

\subsection*{练习5.31（难度7·极难）}

\textbf{L2 知识点梳理：}
这道题综合了定义域、奇偶性、对数不等式和代数运算。
(1) 定义域：真数 \(> 0\)。
(2) 奇偶性：计算 \(f(-x)\)。
(3) 不等式：\(\lg\frac{1-x}{1+x} > 0 = \lg 1\)，利用 \(\lg\) 递增的性质，注意定义域。
(4) 代入 \(f(a)+f(b) = f(c)\)，利用前面的结论。

\textbf{L3 第一步：}
(1) \(\frac{1-x}{1+x} > 0\)，分式不等式。
(2) \(f(-x) = \lg\frac{1+x}{1-x} = \lg\left(\frac{1-x}{1+x}\right)^{-1} = -\lg\frac{1-x}{1+x} = -f(x)\)。
(3) \(\lg\frac{1-x}{1+x} > 0 \Rightarrow \frac{1-x}{1+x} > 1\)，注意定义域。
(4) 利用练习5.25的结论：\(f(a)+f(b) = f\left(\frac{a+b}{1+ab}\right)\)。

\textbf{L4 完整思路：}
(1) \(\frac{1-x}{1+x} > 0 \Rightarrow (1-x)(1+x) > 0 \Rightarrow 1 - x^2 > 0 \Rightarrow x^2 < 1 \Rightarrow -1 < x < 1\)。
(2) 已证，奇函数。
(3) \(\frac{1-x}{1+x} > 1 \Rightarrow \frac{1-x}{1+x} - 1 > 0 \Rightarrow \frac{1-x-(1+x)}{1+x} > 0 \Rightarrow \frac{-2x}{1+x} > 0 \Rightarrow \frac{2x}{1+x} < 0\)。
解这个分式不等式得 \(-1 < x < 0\)。结合定义域，解集为 \((-1, 0)\)。
(4) \(f(a)+f(b) = f(c) \Rightarrow f\left(\frac{a+b}{1+ab}\right) = f(c)\)。
因为 \(f\) 在定义域内单调，所以 \(\frac{a+b}{1+ab} = c\)。注意要求 \(\frac{a+b}{1+ab}\) 在定义域 \((-1,1)\) 内。
